\begin{exercise}{133}
Given a cocartesian fibration $\mathcal E\to\mathcal C\times\Delta^1$, construct a functor $\mathcal E_0\to\mathcal E_1$
which commutes with the projections to $\mathcal C$. Here, $\mathcal E_i$ is the pulled-back cocartesian
fibration along the inclusion $\mathcal C\times\{i\}\to\mathcal C\times\Delta^1$. Show that this functor is a morphism
of cocartesian fibrations. Likewise, construct for a cocartesian fibration $\mathcal E\to\mathcal C\times\Delta^2$ a
$2$-simplex in the $\infty$-category $(\mathrm{Cat}_\infty)_{/\mathcal C}$. If you are eager, do this for a general $n$ instead
of $2$.

\begin{proof}
We give the construction for general $n$.
Write $p\colon\mathcal E\to\mathcal C\times\Delta^n$ and
$p_i\colon\mathcal E_i\to\mathcal C$.
By Corollary~3.2.10, $\mathcal E^{\natural}\to(\mathcal C\times\Delta^n)^{\sharp}$
has the marked left-anodyne lifting property.
This property survives pullback, so the same corollary identifies the
induced marking on $\mathcal E_i$ with its $p_i$-cocartesian edges.

Put $K_i=\Delta^{\{i,\ldots,n\}}$, with inclusion $\iota_i\colon K_i\to\Delta^n$.
Let $S_i$ have $r$-simplices the pairs
\[
 \begin{gathered}
  \gamma\colon\Delta^r\to\mathcal C,\qquad
  h\colon(\Delta^r)^{\flat}\times K_i^{\sharp}\to\mathcal E^{\natural},\\
  ph=(\gamma\mathrm{pr}_1,\iota_i\mathrm{pr}_2),\qquad h\text{ marked}.
 \end{gathered}
\]
Evaluation gives $e_i\colon S_i\to\mathcal E_i$, and tail restriction gives
$R_{ij}\colon S_i\to S_j$, with $R_{jk}R_{ij}=R_{ik}$ and
$p_je_jR_{ij}=p_ie_i$.
A boundary lifting problem for $e_i$ is adjoint to the pushout product
of $\{i\}\to K_i^{\sharp}$ with
$(\partial\Delta^r)^{\flat}\to(\Delta^r)^{\flat}$.
The first is marked left-anodyne by Lemma~3.2.17, since $i$ is the initial
vertex of $K_i$; the pushout product is marked left-anodyne by Lemma~3.2.12.
Thus $e_i$ is a trivial fibration and $S_i$ is an $\infty$-category.

Apply Lemma~3.2.12 also with $\varnothing\to\mathcal E_i^{\natural}$.
The lifting property then extends the inclusion at $i$ to a marked map
\[
 \widehat s_i\colon\mathcal E_i^{\natural}\times K_i^{\sharp}
 \longrightarrow\mathcal E^{\natural}
\]
over $(p_i\mathrm{pr}_1,\iota_i\mathrm{pr}_2)$.
It defines a section $s_i\colon\mathcal E_i\to S_i$ of $e_i$.
Set $F_{ij}=e_jR_{ij}s_i$.
Then $p_jF_{ij}=p_i$ strictly, and $F_{ij}$ preserves cocartesian edges:
the restriction of $\widehat s_i$ at $j$ preserves the induced markings.
Its restriction along each $\{z\}\times\Delta^{\{i,j\}}$ is also
$p$-cocartesian.

Lift the endpoint maps $\mathrm{id}_{S_i}$ and $s_ie_i$ through $e_i$
over $S_i\times J$, with base map $e_i\mathrm{pr}_1$.
Exercises~\ref{ex:59} and~\ref{ex:61} give a natural equivalence
\[
 \begin{gathered}
 \theta_i\colon\Delta^1\to\mathrm{Fun}(S_i,S_i)^{\simeq},\\
 \theta_i(0)=\mathrm{id},\quad\theta_i(1)=s_ie_i,\quad
 e_i\theta_i=\mathrm{const}_{e_i}.
 \end{gathered}
\]
For $i<j$, use enriched composition to define
\[
 \begin{gathered}
 \Phi_{ij}\colon(\Delta^1)^{j-i-1}\to
       \mathrm{Fun}(\mathcal E_i,\mathcal E_j)^{\simeq},\\
 \Phi_{ij}(\mathbf t)=
 e_jR_{j-1,j}\theta_{j-1}(t_{j-1})\cdots R_{i,i+1}s_i.
 \end{gathered}
\]
Here one inserts $\theta_k(t_k)$ at every intermediate $S_k$;
for $j=i+1$ there are no such factors.
At all coordinates zero the value is $F_{ij}$.
Setting $t_k=1$ inserts $s_ke_k$ and splits the product, giving
\[
 \Phi_{ij}|_{t_k=1}=\Phi_{kj}\Phi_{ik},\qquad
 p_j\Phi_{ij}=\mathrm{const}_{p_i}.
\]
The second equality follows from $p_je_jR_{ij}=p_ie_i$ and the strict
equation for $e_i\theta_i$.

In the coherent-nerve model of $\mathrm{Cat}_\infty$
(Definitions~1.2.63 and~2.1.13), send the objects $0,\ldots,n,n+1$
to $\mathcal E_0,\ldots,\mathcal E_n,\mathcal C$.
Send the hom-cubes for $i<j\le n$ to $\Phi_{ij}$, those from $i$ to $n+1$
to the constant map at $p_i$, and identities to identities.
Composition in $\mathfrak C[\Delta^{n+1}]$ inserts the intermediate
coordinate $1$, by the subset-union rule of Observation~1.2.55.
The displayed identities therefore verify every composition, including
those with final object $\mathcal C$.
We obtain an actual $(n+1)$-simplex with last vertex $\mathcal C$,
hence an $n$-simplex of $(\mathrm{Cat}_\infty)_{/\mathcal C}$ by
Corollary~1.4.17.
For $n=1$ its arrow is $F_{01}$; for $n=2$ its comparison is
$\Phi_{02}\colon F_{02}\Rightarrow F_{12}F_{01}$, giving all the requested data.
\end{proof}
\end{exercise}
